% ============================================================================
%  workdev-cheatsheet.tex
%  Quick reference for the workdev tmux workspace command.
%  Requires linuxnotes.sty in the same folder.
%  Compile twice with pdflatex so "Page X of Y" resolves.
% ============================================================================
\documentclass[11pt]{article}

\usepackage[pagenumbers]{linuxnotes}

\linuxnotesmeta{workdev Cheat Sheet}

% Header row style shared by every table: shaded and bold, so it is easy to
% mark as header cells during the Acrobat tagging pass.
\newcommand{\headrow}{\rowcolor{navy!8}}

\begin{document}

\linuxnotestitle{WORK-DEV}{workdev Cheat Sheet}{One command to open a ready-made project workspace}

\noindent
Run \cmd{workdev} \var{NAME} to open a project. Running it again reattaches to the same
workspace, which keeps running in the background until you close it.

\section{Commands}

\par\noindent\begin{minipage}{\linewidth}\centering
\captionof{table}{Every workdev command and option.}
\renewcommand{\arraystretch}{1.15}
\begin{tabular}{>{\raggedright\arraybackslash}p{0.36\linewidth} >{\raggedright\arraybackslash}p{0.56\linewidth}}
\toprule
\headrow
\bfseries Command & \bfseries What it does \\
\midrule
\cmd{workdev} & Pick a project from \cmd{$WORK_DIR} (uses fzf if installed) \\
\cmd{workdev} \var{NAME} & Open the project folder \cmd{$WORK_DIR/}\var{NAME} \\
\cmd{workdev} \var{PATH} & Open any folder by path, such as \cmd{workdev .} \\
\cmd{workdev} \var{NAME} \var{BRANCH} & Switch to git branch \var{BRANCH}, then open the project \\
\cmd{workdev -d} \var{NAME} [\var{BRANCH}] & Build the workspace without attaching \\
\cmd{workdev -l} & List running workspaces \\
\cmd{workdev -k} \var{NAME} & Close the workspace \var{NAME} completely \\
\cmd{workdev -h} & Show the built-in help \\
\bottomrule
\end{tabular}
\end{minipage}\par\medskip

\section{Layout}

The \cmd{dev} window (window 0) has three labeled panes. A \cmd{server} window
(window 1) appears only when \file{.workrc} sets \cmd{SERVER_CMD}.

\par\noindent\begin{minipage}{\linewidth}\centering
\captionof{table}{Panes in the dev window.}
\renewcommand{\arraystretch}{1.15}
\begin{tabular}{>{\raggedright\arraybackslash}p{0.18\linewidth} >{\raggedright\arraybackslash}p{0.18\linewidth} >{\raggedright\arraybackslash}p{0.54\linewidth}}
\toprule
\headrow
\bfseries Pane & \bfseries Position & \bfseries What it shows \\
\midrule
editor [\var{branch}] & Left & Your editor; the label shows the current git branch \\
shell & Top right & A plain terminal in the project folder \\
status & Bottom right & Live \cmd{git status} every 5 seconds, or a file listing outside git \\
\bottomrule
\end{tabular}
\end{minipage}\par\medskip

\section{Settings}

Add these lines to \file{~/.bashrc}, then run \cmd{source ~/.bashrc} to reload it.

\begin{lstlisting}
export PATH="$HOME/.local/bin:$PATH"
export WORK_DIR="$HOME/Projects"
\end{lstlisting}

Optional per-project settings go in \file{.workrc} in the project folder, one
\cmd{KEY=value} per line. The file is read, never executed.

\par\noindent\begin{minipage}{\linewidth}\centering
\captionof{table}{Settings recognized in a .workrc file.}
\renewcommand{\arraystretch}{1.15}
\begin{tabular}{>{\raggedright\arraybackslash}p{0.16\linewidth} >{\raggedright\arraybackslash}p{0.38\linewidth} >{\raggedright\arraybackslash}p{0.36\linewidth}}
\toprule
\headrow
\bfseries Key & \bfseries What it controls & \bfseries Example \\
\midrule
\cmd{EDITOR_CMD} & Editor pane (default: \cmd{$EDITOR .}) & \cmd{EDITOR_CMD=nvim .} \\
\cmd{STATUS_CMD} & Status pane (default: git status) & \cmd{STATUS_CMD=git log --oneline -15} \\
\cmd{SERVER_CMD} & Adds a \cmd{server} window running it & \cmd{SERVER_CMD=python3 -m http.server} \\
\bottomrule
\end{tabular}
\end{minipage}\par\medskip

\section{Common fixes}

\par\noindent\begin{minipage}{\linewidth}\centering
\captionof{table}{Problems and how to fix them.}
\renewcommand{\arraystretch}{1.15}
\begin{tabular}{>{\raggedright\arraybackslash}p{0.27\linewidth} >{\raggedright\arraybackslash}p{0.27\linewidth} >{\raggedright\arraybackslash}p{0.36\linewidth}}
\toprule
\headrow
\bfseries Problem & \bfseries Cause & \bfseries Fix \\
\midrule
\cmd{workdev: command not found} & \file{~/.local/bin} is not on PATH, or the file is not executable & Add the PATH line above, then run \cmd{chmod +x ~/.local/bin/workdev} \\
``Cannot find a project'' & A tilde inside quotes is not expanded & Use \cmd{"$HOME/Projects"}, not \cmd{"~/Projects"} \\
Only some panes appear & Panes were closed with \cmd{exit} & Run \cmd{workdev -k} \var{NAME}, then \cmd{workdev} \var{NAME} \\
``Not switching because of uncommitted changes'' & Edits would carry over to the other branch & Run \cmd{git commit -am "message"} or \cmd{git stash}, then retry \\
``There is no branch named \ldots'' & A typo, or a branch not yet created & Pick from the list shown, or run \cmd{git switch -c} \var{NAME} \\
Editor shows old text after a switch & Files were loaded before the switch & In vim, type \cmd{:e!} and press \key{Enter} \\
\bottomrule
\end{tabular}
\end{minipage}\par\medskip

\warnbox{Remember.}{To step away without losing anything, detach with \key{Ctrl+b} then
\key{d} instead of exiting the panes.}

\vspace{1em}
\noindent{\footnotesize\color{navysoft}\copyright\ 2026 Jason Powell. Part of Work-Dev. Licensed under
\href{https://creativecommons.org/licenses/by/4.0/}{Creative Commons Attribution 4.0 (CC BY 4.0)}.}

\end{document}
